\documentclass[11pt,a4paper]{article}
\usepackage[utf8]{inputenc}
\usepackage[T1]{fontenc}
\usepackage[italian]{babel}
\usepackage{geometry}
\usepackage{hyperref}
\usepackage{tikz}
\usetikzlibrary{automata, positioning, arrows.meta}

\usepackage{tikz-uml}
\usepackage{xcolor}
\usepackage{amsmath,amssymb}
%\setcounter{secnumdepth}{-\maxdimen} % remove section numbering
\usepackage{iftex}
\usepackage{etoolbox}
\usepackage{caption}
\usepackage{subcaption}
\usepackage{graphicx}
\usepackage{float}
\usepackage[none]{hyphenat}
\graphicspath{{figures/}}
\geometry{margin=2.5cm}

\begin{document}
\title{RMI-Chat\\
    \Large{Report di progetto per il corso "distributed algorithms"}
} % Title

\author{\textsc{Alessandro Appio} \\
    \emph{Matricola n. 207228} \\
    corso di laurea Magistrale in Informatica\\
    \emph{296403@studenti.unimore.it}
  }

\date{\today}

\maketitle % Inserts the title, author and other information
\newpage
%----------------------------------------------------------------------------------------
%	Main Part
%----------------------------------------------------------------------------------------
\tableofcontents
\newpage
\maketitle

\section{Introduzione}
Il progetto \texttt{RMI-Chat} realizza un sistema di messaggistica distribuito basato su architettura client-server. I client non comunicano direttamente tra loro, ma attraverso un server centrale che coordina sessioni, canali pubblici e chat private.

\section{Requisiti del sistema}
\subsection{Requisiti funzionali}
\begin{itemize}
  \item \textbf{Connessione client}: il sistema deve consentire al client di connettersi al server tramite URI remoto.
  \item \textbf{Identità utente}: il server deve garantire unicità dello username durante una sessione attiva.
  \item \textbf{Lista utenti connessi}: il client deve poter richiedere la lista degli utenti online.
  \item \textbf{Gestione canali pubblici}: il client deve poter entrare/uscire da un canale pubblico e scambiare messaggi con i membri del canale.
  \item \textbf{Chat privata}: il client deve poter avviare una conversazione privata con un altro utente connesso.
  \item \textbf{Notifiche}: il server deve inviare notifiche asincrone ai client (es. richiesta chat privata).
  \item \textbf{Gestione stato client}: il client deve operare secondo stati espliciti (disconnesso, lobby, canale, privato, uscita).
\end{itemize}

\subsection{Requisiti non funzionali}
\begin{itemize}
  \item \textbf{Modularità}: separazione tra componenti client, server e moduli comuni.
  \item \textbf{Configurabilità}: parametri server esterni al codice (\texttt{config.toml}).
  \item \textbf{Usabilità}: interfaccia testuale guidata da menu.
  \item \textbf{Estendibilità}: possibilità di aggiungere nuove operazioni remote senza modificare l'interfaccia utente di base.
\end{itemize}

\section{Tecnologie utilizzate}
Il progetto è stato sviluppato in Python, utilizzando:
\begin{itemize}
  \item \textbf{Pyro5} per l'RMI (invocazioni remote tra processi);
  \item \textbf{inquirer} per l'interfaccia testuale interattiva lato client;
  \item file \textbf{TOML} per la configurazione del server (host, nome, canali, messaggio di benvenuto).
\end{itemize}

\section{Architettura del sistema}
L'architettura è di tipo client-server centralizzato:
\begin{itemize}
  \item il \textbf{server} espone i metodi remoti per connessione utenti, gestione canali e inoltro messaggi;
  \item ogni \textbf{client} si collega al server e registra un proprio oggetto remoto per ricevere messaggi e notifiche.
\end{itemize}

\subsection{Diagramma architetturale}

\begin{figure}[h!]
    \centering
    \begin{tikzpicture}[->, >=Stealth, auto, node distance=5cm, semithick]
        \tikzstyle{state}=[circle, draw, minimum size=5mm]

        \node[state] (S) {Server};
        \node[state, above left of=S] (A) {Client A};
        \node[state, left of=S] (B) {Client B};
        \node[state, below left of=S] (C) {Client C};

        \path
        (S) edge node {} (A)
        (A) edge node {RPC/Callback} (S)
        (S) edge node {} (B)
        (B) edge node {RPC/Callback} (S)
        (S) edge node {} (C)
        (C) edge node {RPC/Callback} (S);
    \end{tikzpicture}
    \caption{Connection diagram}
    \label{fig:placeholder}
\end{figure}

Il server mantiene lo stato globale (utenti, canali, richieste private), mentre i client mantengono solo lo stato locale della sessione utente.

\section{Implementazione dei componenti}
\subsection{Lato server}
Il server carica la configurazione da \texttt{config/config.toml}, avvia un daemon Pyro e registra l'oggetto \texttt{Chat}.

\noindent Le funzionalità principali implementate sono:
\begin{itemize}
  \item connessione/disconnessione utenti con controllo dell'unicità username;
  \item visualizzazione utenti connessi;
  \item ingresso/uscita da canali pubblici;
  \item richiesta e gestione di chat private;
  \item inoltro messaggi nei canali e nelle chat private;
  \item notifiche asincrone ai client.
\end{itemize}

%
%  CLASS UML DIAGRAM HERE
%
\begin{figure}[H]
    \centering
    \begin{tikzpicture}
      % 1. Define Class 'Person'
      \umlclass[x=0, y=0]{Chat}{
        - clients : ClientInfo[] \\
        - channels : String[]
      }{
        + connect() : int \\
        + leave() \\
        + enterChannel() \\
        + leaveChannel() \\
        + enterPrivate() : int \\
        + leavePrivate()
      }

      % 2. Define Class 'Student' (Positioned right of Person)
      \umlclass[x=7, y=0]{ClientInfo}{
        - uri : String\\
        - username : String\\
        - channel : String\\
        - privateRelative : String\\
        - inChannel : bool\\
        - inPrivate : bool
      }{
        + createProxy() : Pyro5.api.Proxy
      }
    \end{tikzpicture}
    \caption{Classi per il server}
    \label{fig:placeholder}
\end{figure}


\subsection{Lato client}
Il client avvia un'interfaccia CLI e una macchina a stati che regola il comportamento dell'applicazione.\\
Gli stati implementati sono:
\begin{itemize}
  \item \texttt{DISCONNECTED}
  \item \texttt{LOBBY}
  \item \texttt{CHANNEL}
  \item \texttt{PRIVATE}
  \item \texttt{QUIT}
\end{itemize}

In base allo stato, l'utente può eseguire operazioni diverse (connessione, scelta canale, avvio chat privata, invio messaggi, uscita).

L'intera logica della macchina a stati è gestita nella classe StateMachine e gli stati sono definiti in una classe che funge da enum chiamata State.

Ogni Client espone un oggetto con pyro con tre metodi:
\begin{itemize}
    \item \textbf{receive}: per la ricezione di messaggi da altri client quando si trova nello stato CHANNEL o PRIVATE
    \item \textbf{notify}: per la ricezione di messaggi dal server in qualsiasi stato
    \item \textbf{quitPrivate}: per comunicare al client in una chat privata che la chat si è conclusa
\end{itemize}

\begin{figure}[H]
    \centering
    \begin{tikzpicture}
      % 1. Define Class 'Person'
      \umlclass[x=0, y=0]{Client}{
        - state: State\\
        - username: String\\
        - uri: String\\
        - chat: Pyro5.api.Proxy\\
        - currentChannel: String\\
        - otherPrivate: String
      }{
        + receive(sender, msg) \\
        + notify(msg) \\
        + quitPrivate(msg)
      }

      % 2. Define Class 'Student' (Positioned right of Person)
      \umlclass[x=7, y=1]{StateMachine}{
        - client: Client
      }{
        + disconnected()\\
        + lobby()\\
        + channel()\\
        + private()
      }

      \umlclass[x=7, y=-3]{State}{
        - DISCONNECTED \\
        - LOBBY \\
        - CHANNEL \\
        - PRIVATE \\
      }{}
    \end{tikzpicture}
    \caption{Classi per il client}
    \label{fig:placeholder}
\end{figure}

\section{Protocolli di comunicazione}
\subsection{Protocollo client-server adottato}
Il protocollo è di tipo \textbf{client-server} su chiamate remote sincrone (RPC) con callback asincrone verso i client. Non è adottato un protocollo peer-to-peer.

\subsection{Diagramma di sequenza -- connessione}
\begin{figure}[H]
    \centering
    \begin{tikzpicture}[
        node distance=2.5cm,
        state/.style={rectangle, draw, rounded corners, minimum height=0.8cm, align=center, fill=gray!10},
        arrow/.style={->, >=stealth, thick}
    ]
        % Nodi
        \node[state] (disc) {Disconnesso\\Menu disconnesso};
        \node[state, right=of disc] (check) {Connessione\\Esito};
        \node[state, right=of check] (welc) {Benvenuto};
        \node[state, right=of welc] (conn) {Connesso\\Menu lobby};

        % Transizioni
        \draw[arrow] (disc) -- node[above, font=\small]{\texttt{connect()}} (check);
        \draw[arrow] (check) -- node[above, font=\small]{Esito = 0} node[below, font=\small]{\texttt{welcome()}} (welc);
        \draw[arrow] (welc) -- node[above, font=\small] {} (conn);
        \draw[arrow] (check.south) -- ++(0,-0.5) -| node[near start, below, font=\small]{Esito = Errore} (disc.south);

    \end{tikzpicture}
    \caption{Macchina a stati della connessione lato Client.}
\end{figure}

\subsection{Diagramma di sequenza -- messaggio in canale}
\begin{figure}[h]
    \centering
    \begin{tikzpicture}[
        node distance=2cm and 3cm,
        state/.style={rectangle, draw, rounded corners, minimum height=0.9cm, align=center, fill=gray!10},
        arrow/.style={->, >=stealth, thick}
    ]

        % Nodi
        \node[state] (enter) {Ingresso\\Canale};
        \node[state, below=of enter] (idle) {In Attesa\\Messaggi};
        \node[state, right=of idle] (id_chan) {Identificazione\\Canale};
        \node[state, right=of id_chan] (loop) {Inoltro Broadcast\\(per ogni client)};

        % Transizioni
        \draw[arrow] (enter) -- node[right, font=\small]{\texttt{enterChannel(uri, channel)}} (idle);
        \draw[arrow] (idle) -- node[above, font=\small]{\texttt{send(uri, msg)}} (id_chan);
        \draw[arrow] (id_chan) -- node[above, font=\small]{Canale} node[below, font=\small]{Identificato} (loop);
        \draw[arrow] (loop) edge[loop above] node[font=\small, align=center]{\texttt{receive()}\\a client $D_i$} (loop);
        \draw[arrow] (loop.south) -- ++(0,-0.6) -| node[near start, below, font=\small]{Fine Lista Client} (idle.south);

    \end{tikzpicture}
    \caption{Macchina a stati dell'inoltro messaggio lato Server.}
\end{figure}

\subsection{Diagramma di sequenza -- avvio chat privata}

Prendiamo in considerazione due client, chiamati $A$ e $B$, di seguito il diagramma che riguarda $A$.

\begin{figure}[h]
    \centering
    \begin{tikzpicture}[
        node distance=1.5cm and 3cm,
        state/.style={rectangle, draw, rounded corners, minimum height=0.9cm, align=center, fill=gray!10},
        arrow/.style={->, >=stealth, thick}
    ]

        % Nodi
        \node[state] (enter) {Richiesta chat\\privata};
        \node[state, below=of enter] (wait) {In Attesa\\agreement di $B$};
        \node[state, below=of wait] (idle) {Ingresso chat};
        \node[state, right=of idle] (id_chan) {Identificazione\\uri di $B$};
        \node[state, right=of id_chan] (loop) {Inoltro a $B$};

        % Transizioni
        \draw[arrow] (enter) -- node[right, font=\small]{\texttt{enterChannel(uri, channel)}} (wait);
        \draw[arrow] (wait) -- node[right, font=\small]{Agreement} (idle);
        \draw[arrow] (idle) -- node[above, font=\small]{\texttt{send(uri, msg)}} (id_chan);
        \draw[arrow] (id_chan) -- node[above, font=\small]{uri} node[below, font=\small]{Identificato} (loop);
        \draw[arrow] (loop) edge[loop above] node[font=\small, align=center]{\texttt{receive()}\\a client $B$} (loop);
        \draw[arrow] (loop.south) -- ++(0,-0.6) -| node[near start, below, font=\small]{} (idle.south);

    \end{tikzpicture}
    \caption{Macchina a stati dell'inoltro messaggio lato Server.}
\end{figure}

\section{Funzionamento operativo}
Flusso operativo tipico:
\begin{enumerate}
  \item avvio del server;
  \item avvio di uno o più client;
  \item connessione client al server con username;
  \item ingresso in canale pubblico o apertura chat privata;
  \item scambio di messaggi in tempo reale;
  \item uscita dalla conversazione e disconnessione.
\end{enumerate}

\section{Risultati ottenuti}
Il sistema realizzato soddisfa i requisiti principali di una chat distribuita didattica:
\begin{itemize}
  \item comunicazione remota funzionante tra nodi distinti;
  \item supporto a conversazioni sia pubbliche sia private;
  \item gestione dello stato client in modo chiaro;
  \item configurabilità del server senza modifiche al codice.
\end{itemize}

\section{Conclusioni}
Il progetto ha permesso di applicare in modo pratico i concetti di sistemi distribuiti e invocazione remota di metodi. La soluzione implementata è modulare, estendibile e adatta a dimostrare le principali dinamiche di coordinamento client-server in un contesto di messaggistica.

\end{document}
